% BEGIN REV09_PELL_MOMENTO_PREVIO_IV
\subsection{Oriented depth moment}
\label{corpus9:iv:pell-momento}

The realization of the depth register with a quarter-turn
defines the parametric moment
\[
 \mathcal C_2(r)=
   \sum_{k\geq0}\frac{(-1)^kr^{2k+1}}{(2k+1)^2},
 \qquad 0<r<1,\qquad
 \mathcal D=r\frac{d}{dr}.
\]
Local uniform convergence permits termwise differentiation:
\[
 \mathcal D\mathcal C_2(r)=\arctan r,\qquad
 \mathcal D^2\mathcal C_2(r)=\frac{r}{1+r^2}.
\]
Indeed, the first differentiated series is the integral from zero
of the geometric series \(1/(1+r^2)\). A second logarithmic
differentiation gives the last equality. The conditions
\(\mathcal C_2(0)=\mathcal D\mathcal C_2(0)=0\) fix the
integration constants and give
\[
 \mathcal C_2(r)=\int_0^r\frac{\arctan t}{t}\,dt
   =\int_0^r\frac{\log(r/t)}{1+t^2}\,dt.
\]
Domination by \(\sum_{n\geq1}n^{-2}\) retains the limit
\(\mathcal C_2(1)=G_{\mathrm{Cat}}\). The second derivative
at \(r=1\) is understood through its Abel limit.
% END REV09_PELL_MOMENTO_PREVIO_IV
